\section{Phát hiện từ lặp liên tiếp [Cơ bản]}
\subsection{Phân tích \& Thiết kế (M0)}

\begin{itemize}

    \item \textbf{Input:}

    File \texttt{text.txt} chứa nhiều dòng văn bản ASCII.

    \item \textbf{Xử lý:}

    \begin{itemize}

        \item \textbf{Case biên:}

        \begin{itemize}
            \item Dòng rỗng $\rightarrow$ bỏ qua
            \item Dòng chỉ có 1 token $\rightarrow$ không có lặp
            \item Không có từ lặp liên tiếp
            $\rightarrow$
            file output rỗng
        \end{itemize}

        \item \textbf{Module 1 (Đọc file):}

        \begin{itemize}
            \item Hàm:
\begin{lstlisting}
readLines(path, vector<string>&)
\end{lstlisting}

            \item Kiểm tra mở file
        \end{itemize}

        \item \textbf{Module 2 (Normalize chuỗi):}

        \begin{itemize}
            \item Hàm:
\begin{lstlisting}
normalizeLine(const string& s)
\end{lstlisting}

            \item Chuyển lowercase
            \item Ký tự không phải chữ/số
            $\rightarrow$
            space
            \item Gom nhiều space thành 1
        \end{itemize}

        \item \textbf{Module 3 (Tokenize):}

        \begin{itemize}
            \item Hàm:
\begin{lstlisting}
vector<string> tokenize(...)
\end{lstlisting}
        \end{itemize}

        \item \textbf{Module 4 (Phát hiện lặp liên tiếp):}

        \begin{itemize}
            \item Duyệt token theo từng dòng
            \item Nếu:
\begin{lstlisting}
tokens[i] == tokens[i-1]
\end{lstlisting}

            $\rightarrow$
            tăng count

            \item Nếu count $\ge 2$
            $\rightarrow$
            ghi kết quả
        \end{itemize}

        \item \textbf{Module 5 (Ghi output):}

        \begin{itemize}
            \item File:
\begin{lstlisting}
repeats.txt
\end{lstlisting}

            \item Format:
\begin{lstlisting}
line=<i> word=<w> count=<k>
\end{lstlisting}

            \item \texttt{i} là số dòng bắt đầu từ 1
        \end{itemize}

    \end{itemize}

    \item \textbf{Output:}

\begin{lstlisting}
line=<i> word=<w> count=<k>
\end{lstlisting}

\end{itemize}

\subsection{Mã nguồn C++11 (Tách module)}

\begin{lstlisting}
#include <iostream>
#include <vector>
#include <fstream>
#include <string>
#include <sstream>
#include <cctype>

using namespace std;

struct RepeatInfo {
    int line;
    string word;
    int count;
};

bool readLines(const string& path,
               vector<string>& lines) {

    ifstream fin(path);

    if (!fin) return false;

    string s;

    while (getline(fin, s)) {

        lines.push_back(s);
    }

    return true;
}

string normalizeLine(const string& s) {

    string result;

    bool lastSpace = true;

    for (char c : s) {

        if (isalnum((unsigned char)c)) {

            result +=
                (char)tolower((unsigned char)c);

            lastSpace = false;
        }
        else {

            if (!lastSpace) {

                result += ' ';
                lastSpace = true;
            }
        }
    }

    if (!result.empty() &&
        result.back() == ' ') {

        result.pop_back();
    }

    return result;
}

vector<string> tokenize(
    const string& normalized) {

    vector<string> tokens;

    stringstream ss(normalized);

    string token;

    while (ss >> token) {

        tokens.push_back(token);
    }

    return tokens;
}

bool writeRepeats(
    const string& path,
    const vector<RepeatInfo>& repeats) {

    ofstream fout(path);

    if (!fout) return false;

    for (const RepeatInfo& r : repeats) {

        fout << "line="
             << r.line
             << " word="
             << r.word
             << " count="
             << r.count
             << "\n";
    }

    return true;
}

int main() {

    vector<string> lines;

    if (!readLines("text.txt", lines)) {

        cerr << "Loi mo file input!\n";
        return 1;
    }

    vector<RepeatInfo> repeats;

    for (size_t lineIdx = 0;
         lineIdx < lines.size();
         ++lineIdx) {

        string normalized =
            normalizeLine(lines[lineIdx]);

        vector<string> tokens =
            tokenize(normalized);

        if (tokens.empty()) continue;

        string current = tokens[0];
        int count = 1;

        for (size_t i = 1;
             i < tokens.size();
             ++i) {

            if (tokens[i] == current) {

                count++;
            }
            else {

                if (count >= 2) {

                    repeats.push_back({
                        (int)lineIdx + 1,
                        current,
                        count
                    });
                }

                current = tokens[i];
                count = 1;
            }
        }

        if (count >= 2) {

            repeats.push_back({
                (int)lineIdx + 1,
                current,
                count
            });
        }
    }

    if (!writeRepeats("repeats.txt",
                      repeats)) {

        cerr << "Loi mo file output!\n";
        return 1;
    }

    return 0;
}
\end{lstlisting}

\subsection{Kế hoạch kiểm thử (Test Cases)}

\textbf{Test Case 1 (Thông thường)}

\textbf{Input (text.txt):}

\begin{lstlisting}
hello hello world
test test test data
a b c
\end{lstlisting}

\textbf{Expected Output (repeats.txt):}

\begin{lstlisting}
line=1 word=hello count=2
line=2 word=test count=3
\end{lstlisting}

\vspace{0.5cm}

\textbf{Test Case 2 (Không có lặp)}

\textbf{Input (text.txt):}

\begin{lstlisting}
one two three
apple banana orange
\end{lstlisting}

\textbf{Expected Output (repeats.txt):}

\begin{lstlisting}

\end{lstlisting}

\vspace{0.5cm}

\textbf{Test Case 3 (Case biên -- dòng rỗng)}

\textbf{Input (text.txt):}

\begin{lstlisting}

hello hello

wow wow wow
\end{lstlisting}

\textbf{Expected Output (repeats.txt):}

\begin{lstlisting}
line=2 word=hello count=2
line=4 word=wow count=3
\end{lstlisting}